\section{Compact fillings and Euclidean current charts}
\label{sec:currents}

We first reduce the filling problem to a compact ambient space. Standard
compactness and filling continuity then turn a hypothetical failure of the
sharp bound into a normalized extremizing cycle. To vary this cycle
without losing the sharp coefficient, we derive exact Euclidean chart
mass and density formulas. These local formulas apply to every integral
current. All current dimensions in this section are finite positive
integers; the ambient space has no dimension bound.

\subsection{Reduction to a compact convex ball}

\begin{lemma}[Compact reduction]
\label{lem:compact-reduction}
Suppose that the asserted filling inequality in a fixed cycle dimension
holds in every compact \(\operatorname{CAT}(0)\) space. Then it holds for
every compactly supported integral cycle of that dimension in every proper
\(\operatorname{CAT}(0)\) space, with a compactly supported integral filling.
\end{lemma}

\begin{proof}
Let \(T\) be the cycle in a proper \(\operatorname{CAT}(0)\) space \(X\).
If \(T=0\), take the zero filling. Otherwise choose a closed ball
\(Y=\overline B(o,R)\) whose interior contains \(\spt T\).
Properness makes \(Y\) compact. Triangle comparison implies that balls
are convex, so the induced metric makes \(Y\) a compact
\(\operatorname{CAT}(0)\) space.

For \(x\in X\), let \(P(x)\) be the point of the segment \([o,x]\)
at distance \(\min\{R,d(o,x)\}\) from \(o\).
Geodesics in a \(\operatorname{CAT}(0)\) space are unique. Comparing two
such radial points with their counterparts in the Euclidean comparison
triangle, and then using nonexpansion of Euclidean radial projection onto
a closed ball, gives
\[
 d(P(x),P(x'))\le d(x,x').
\]
One can also identify \(P(x)\) as the nearest point of \(Y\): the triangle
inequality bounds the distance from \(x\) to \(Y\) below by
\(\max\{0,d(o,x)-R\}\), and the radial point attains this bound.

Lipschitz pushforward therefore gives an integral cycle \(P_\#T\) in
\(Y\) with \(\M(P_\#T)\le\M(T)\). If
\(\iota:Y\hookrightarrow X\) denotes inclusion, locality gives
\(\iota_\#P_\#T=T\), because \(\iota\circ P\) is the identity on a
neighborhood of \(\spt T\). Apply the compact-space inequality to
\(P_\#T\) and include its filling into \(X\). Its mass does not increase,
its boundary is \(T\), and its support is contained in the compact set
\(Y\).
\end{proof}

We work henceforth in a fixed compact \(\operatorname{CAT}(0)\) space
\(Y\). The elementary comparison facts just used, and the midpoint and
geodesic contraction inequalities used below, follow from the defining
triangle comparison; see also \cite[Chapter~II.2]{BridsonHaefliger1999}.

For an integral \(k\)-cycle \(B\) in \(Y\), define
\begin{equation}
\label{eq:filling-volume-definition}
 V(B)=\inf\{\M(S):S\in\I_{k+1}(Y),\ \partial S=B\}.
\end{equation}
The mass measure of a current \(C\) is denoted by \(\|C\|\), and its total
mass by \(\M(C)\). We use the Ambrosio--Kirchheim convention for the action
\(C(b,\pi_1,\ldots,\pi_k)\): the first coefficient is bounded Lipschitz,
and the differentiated scalar coordinates are Lipschitz. The mass bound
extends the first coefficient to bounded Borel functions. In particular,
if a finite Borel measure \(\eta\) satisfies
\begin{equation}
\label{eq:mass-controlling-measure}
 |C(b,\pi_1,\ldots,\pi_k)|\le\int |b|\dd\eta
\end{equation}
for every bounded Lipschitz \(b\) and every tuple of \(1\)-Lipschitz
coordinates, then \(\|C\|\le\eta\). This is the minimality property in
the definition of mass.

\subsection{The current-theoretic inputs}

The following theorem records the standard foundations in the precise
form used here. The current representation, compactness, closure and
slicing assertions are theorems of Ambrosio--Kirchheim; the product and
filling assertions use Wenger's results. We include the verifications
specific to the present ambient space.

\begin{theorem}[Current background in a compact \(\operatorname{CAT}(0)\) space]
\label{thm:current-background}
Let \(Y\) be compact and \(\operatorname{CAT}(0)\), and let \(k\ge1\).
The following statements hold.
\begin{enumerate}[label=\textup{(\roman*)}]
\item Lipschitz pushforward preserves integral currents, commutes with
boundary, and satisfies the Lipschitz mass bound. Restrictions by bounded
Borel coefficients have the usual mass-measure bound; for a Borel set
\(E\subset Y\),
\[
 \|C\restr E\|=\|C\|\restr E.
\]
The restricted current need not be normal. If \(u:Y\to\R\)
is Lipschitz and \(C\in\I_k(Y)\), then
\(C\restr\{u>t\}\) is integral for almost every \(t\).
In particular its boundary is an integral \((k-1)\)-cycle.

\item If \(C_j\in\I_k(Y)\) and
\(\sup_j(\M(C_j)+\M(\partial C_j))<\infty\), a subsequence converges
weakly to an integral current. Boundary commutes with weak convergence,
and mass is lower semicontinuous.

\item Every \(C\in\I_k(Y)\) admits a representation by bi-Lipschitz
maps \(\phi_i:E_i\to Y\), where \(E_i\subset\R^k\) are bounded Borel
sets and the images \(Z_i=\phi_i(E_i)\) are pairwise disjoint Borel sets,
and by signed integer multiplicities \(\theta_i\in L^1(E_i)\), such that
\begin{equation}
\label{eq:chart-action}
 C(b,\pi_1,\ldots,\pi_k)
 =\sum_i\int_{E_i}\theta_i(z)b(\phi_i(z))
       \det D\bigl((\pi_1,\ldots,\pi_k)\circ\phi_i\bigr)(z)\dd z.
\end{equation}
The chart summands have summable masses. Zero multiplicities may be
discarded. Scalar functions on a Borel chart domain are differentiated
by Lipschitz extension to \(\R^k\); their derivatives are independent
of that extension at almost every density point of the domain.

\item For \(h>0\), the whole-current product \([0,h]\times C\), with
the Euclidean product metric, is an integral \((k+1)\)-current and obeys
\begin{equation}
\label{eq:interval-product-boundary}
 \partial([0,h]\times C)
 =[h]\times C-[0]\times C-[0,h]\times\partial C.
\end{equation}
Writing \(b_t(x)=b(t,x)\) and \(\pi_{a,t}(x)=\pi_a(t,x)\), its action is
\begin{equation}
\label{eq:interval-product-action}
\begin{split}
 &([0,h]\times C)(b,\pi_1,\ldots,\pi_{k+1})\\
 &\quad=\sum_{a=1}^{k+1}(-1)^{a+1}\int_0^h
 C\bigl(b_t\,\partial_t\pi_a(t,\cdot),
          \pi_{1,t},\ldots,\widehat{\pi_{a,t}},\ldots,
          \pi_{k+1,t}\bigr)\dd t.
\end{split}
\end{equation}
The hat indicates omission. The time derivative is interpreted almost
everywhere and as a bounded Borel first coefficient in the current action.

\item There are finite constants \(K_k,L_k\), independent of the cycle,
such that every integral \(k\)-cycle \(B\) in \(Y\) satisfies
\begin{equation}
\label{eq:coarse-fillings}
 V(B)\le K_k\M(B)^{(k+1)/k},\qquad
 V(B)\le L_k\M(B).
\end{equation}
Every such \(B\) has a mass-minimizing integral filling in \(Y\).
If integral \(k\)-cycles \(B_j\) have bounded masses and converge weakly
to an integral cycle \(B\), then
\begin{equation}
\label{eq:filling-continuity}
 V(B_j-B)\longrightarrow0,
 \qquad V(B_j)\longrightarrow V(B).
\end{equation}
\end{enumerate}
\end{theorem}

\begin{proof}
The pushforward and Borel-coefficient properties are part of the basic
current calculus of \cite{AK2000}. The exact restriction identity also
follows from the restriction mass bound and the decomposition
\(C=C\restr E+C\restr(Y\setminus E)\), using minimality of the mass
measure in both directions. The representation in (iii) is the
integer-rectifiable representation
\cite[Lemma~4.1 and Theorem~4.5]{AK2000}; subdivision makes chart domains
bounded and disjointification makes their images disjoint. The slicing
assertion is \cite[Theorems~5.6--5.7]{AK2000}. The compactness and closure
theorems \cite[Theorems~5.2 and~8.5]{AK2000} give (ii): completeness is
automatic, the bounded normal masses are assumed, and compactness of
\(Y\) supplies uniform tightness. The product statements, including the
action convention, are
\cite[Definition~3.1 and Theorem~3.2]{Wenger2007Flat}. They apply because
an integral current is normal and its support in \(Y\) is bounded.

Here is the geometric verification of the filling hypotheses. Choose a
point \(o\) in the support of a nonzero cycle \(B\), and let
\(H(t,x)\) be the point at fraction \(t\) along \([o,x]\).
Triangle comparison gives
\[
 d(H(t,x),H(t,x'))\le t\,d(x,x'),\qquad
 d(H(t,x),H(s,x))=|t-s|d(o,x).
\]
Thus \(H\) is jointly Lipschitz on \([0,1]\times Y\).
On the support of \(B\), its time speed is bounded by
\(D=\diam(\spt B)\), and its spatial Lipschitz constant is at most one.
In each of the \(k+1\) terms of \eqref{eq:interval-product-action}, a
\(1\)-Lipschitz scalar test after composition with \(H\) has time
derivative at most \(D\), and the remaining spatial derivatives have
Lipschitz bounds at most one. The mass bound therefore gives
\begin{equation}
\label{eq:cone-mass-bound}
 \M\bigl(H_\#([0,1]\times B)\bigr)
 \le (k+1)D\M(B).
\end{equation}
The boundary is \(B\): the terminal map is the identity, the initial
map is constant and hence has zero pushforward in positive dimension,
and \(\partial B=0\). The filling is supported on segments from \(o\)
to \(\spt B\), all lying within distance \(D\) of \(o\).
Consequently the space has local cone inequalities in every positive
current dimension. Taking \(D\le\diam Y\) gives the second estimate
in \eqref{eq:coarse-fillings}, for example with
\(L_k=(k+1)\diam Y\).

The complete space \(Y\) has the convex geodesic bicombing given by its
unique geodesics. Wenger's Euclidean-exponent isoperimetric theorem
\cite[Theorem~1.2 and Corollary~1.4]{Wenger2005Iso}, or its cone-inequality
form applied successively in dimension, gives the first estimate in
\eqref{eq:coarse-fillings} for every \(k\ge1\). No sharp value of
\(K_k\) is used here.

For a fixed boundary \(B\), the cone gives at least one filling. A
minimizing sequence has bounded mass and fixed boundary, so (ii) gives
an integral weak limit with that boundary. Lower semicontinuity proves
attainment. Finally, \(Y\) is complete and geodesic, hence quasiconvex,
and it has the local cone inequalities in dimensions \(1,\ldots,k\)
just verified. The cycle assertion in
\cite[Theorem~1.4]{Wenger2007Flat} applies to the weakly convergent
sequence \(B_j\): its normal masses are bounded because its boundary
masses are zero. It gives \(V(B_j-B)\to0\) directly. Finiteness and
subadditivity of filling volume imply
\[
 |V(B_j)-V(B)|\le V(B_j-B),
\]
which proves the second limit.
\end{proof}

\subsection{An extremizing cycle and its normalization}

Fix a cycle dimension \(n\ge2\), and put \(q_0=(n+1)/n\).
The following variational device turns a hypothetical violation of the
sharp filling bound into a cycle with a controlled first variation.

\begin{lemma}[Positive extremizer]
\label{lem:extremizer}
Let \(c>0\). If the functional
\[
 B\longmapsto V(B)-c\M(B)^{q_0}
\]
is positive on some integral \(n\)-cycle in \(Y\), it attains a positive
maximum on such cycles. A maximizing cycle \(A\), with
\(m=\M(A)>0\), satisfies
\begin{equation}
\label{eq:extremizer}
 c\bigl(m^{q_0}-\M(B)^{q_0}\bigr)\le V(A-B)
 \qquad\text{for every integral \(n\)-cycle }B\text{ in }Y.
\end{equation}
\end{lemma}

\begin{proof}
The linear bound in \eqref{eq:coarse-fillings} implies
\[
 V(B)-c\M(B)^{q_0}\le L_n\M(B)-c\M(B)^{q_0}.
\]
The right side has finite supremum, and positivity implies
\(\M(B)<(L_n/c)^n\). A positive maximizing sequence therefore has
uniformly bounded masses and zero boundaries. By
Theorem~\ref{thm:current-background}, a subsequence converges weakly to an
integral cycle \(A\). Filling continuity and lower semicontinuity of
mass give
\[
 V(A)-c\M(A)^{q_0}
 \ge\limsup_j\bigl(V(B_j)-c\M(B_j)^{q_0}\bigr).
\]
Thus \(A\) attains the positive supremum, and in particular
\(\M(A)>0\). Maximality and subadditivity now give
\[
 c\bigl(\M(A)^{q_0}-\M(B)^{q_0}\bigr)
 \le V(A)-V(B)\le V(A-B),
\]
which is \eqref{eq:extremizer}.
\end{proof}

Suppose that the desired sharp inequality fails in dimension \(n\) in
a compact \(\operatorname{CAT}(0)\) space. Attainment in
Theorem~\ref{thm:current-background} makes the failure a strict inequality
\(V(B)>c_n\M(B)^{q_0}\) for some nonzero cycle. Choose
\(c>c_n\) below this ratio, and apply Lemma~\ref{lem:extremizer}.
We write
\(\mu=\|A\|\) and \(m=\M(A)\).

Multiplication of all distances by a positive factor \(a\) multiplies
the mass of a \(j\)-current by exactly \(a^j\). The upper bound follows
from Lipschitz pushforward under the identity map, and the reverse bound
follows by applying the same estimate to its inverse. This identifies
the integral chain groups before and after scaling and gives the same
scaling law for filling volume. Both terms of the maximizing functional
therefore scale by \(a^{n+1}\), so maximality is preserved. Choose the
factor to arrange
\begin{equation}
\label{eq:normalization}
 m=((n+1)c)^{-n}<s_n,
 \qquad cq_0m^{1/n}=\frac1n.
\end{equation}
The strict inequality follows from
\(c>c_n=1/((n+1)s_n^{1/n})\).
We retain the notation \(Y,A,\mu,m\) for these rescaled objects. The
space is still compact and \(\operatorname{CAT}(0)\), and the
current-theoretic inputs remain applicable.

\subsection{Metric differentiation and Euclidean chart norms}

The contradiction now rests on an extremizing cycle of mass $m<s_n$.
To compare its mass loss under radial variation with the mass swept out,
we need chart formulas with their exact Euclidean coefficients. We prove
these formulas for a general integral current $C$: first identify the
metric differential, then recover the mass by scalar current tests, and
finally use integer multiplicity to obtain a lower density.

Fix \(C\in\I_k(Y)\) and charts as in
Theorem~\ref{thm:current-background}. Kirchheim's metric differentiation
theorem \cite[Theorem~2]{Kirchheim1994}, in the two-moving-point formulation
\cite[Definition~(1.3) and Theorem~1.1]{Duda2007}, gives at almost every
chart point \(z\) a seminorm \(N_z\) such that
\begin{equation}
\label{eq:two-point-metric-differential}
 d\bigl(\phi_i(z+w),\phi_i(z+w')\bigr)
 =N_z(w-w')+o(|w|+|w'|)
\end{equation}
as the two domain points tend to \(z\) within \(E_i\).

To apply the Euclidean-domain statement to a Borel domain, first embed
\(Y\) isometrically into \(\ell^\infty(Y)\). More explicitly, choose
\(z_*\in E_i\) and use the coordinates
\[
 z\longmapsto d(\phi_i(z),y)-d(\phi_i(z_*),y),\qquad y\in Y.
\]
Each coordinate is \(L\)-Lipschitz for \(L=\Lip(\phi_i)\) and vanishes
at \(z_*\). By McShane's scalar extension theorem, extend each separately
to \(\R^k\) with the same Lipschitz constant; see
\cite[Theorem~2.1 and Proposition~2.2]{Rieffel2006} for the scalar and
coordinatewise formulations. These extensions have absolute value at most
\(L|z-z_*|\), uniformly over the coordinates, so together they define
an \(\ell^\infty(Y)\)-valued Lipschitz map. Metric differentiation of
this extension gives \eqref{eq:two-point-metric-differential} on the
original domain. This auxiliary extension is used only for differentiation;
all curvature comparisons below concern original image points in \(Y\).

\begin{lemma}[Euclidean chart metrics]
\label{lem:euclidean-chart-metrics}
For almost every density point \(z\in E_i\), the seminorm in
\eqref{eq:two-point-metric-differential} is a Euclidean norm. Thus there
is a measurable positive definite matrix \(G_i(z)\) with
\begin{equation}
\label{eq:chart-metric}
 N_z(w)^2=w^{\mathsf T}G_i(z)w,
 \qquad J_i(z)=\sqrt{\det G_i(z)}.
\end{equation}
In particular,
\begin{equation}
\label{eq:centered-metric-differential}
 d(\phi_i(z+w),\phi_i(z))=|w|_{G_i(z)}+o(|w|)
 \quad (z+w\in E_i).
\end{equation}
\end{lemma}

\begin{proof}
At a density-one point \(z\) of \(E_i\), every fixed finite configuration
in \(\R^k\) can be approximated by points of \((E_i-z)/t\) as
\(t\downarrow0\). Indeed, a point of such a configuration that stayed
a positive distance from \((E_i-z)/t\) along a sequence of scales would
give a missing Euclidean ball of radius comparable to \(t\) at distance
\(O(t)\) from \(z\), contradicting density one. Applying this observation
to a single direction and using the bi-Lipschitz lower bound for
\(\phi_i\) shows that \(N_z(w)\ge\Lip(\phi_i^{-1})^{-1}|w|\).
The upper Lipschitz bound gives \(N_z(w)\le\Lip(\phi_i)|w|\).
Thus \(N_z\) is a norm.

For any four points \(a,b,c,d\) of a \(\operatorname{CAT}(0)\) space,
\begin{equation}
\label{eq:cat0-quadrilateral}
 d(a,c)^2+d(b,d)^2
 \le d(a,b)^2+d(b,c)^2+d(c,d)^2+d(d,a)^2.
\end{equation}
To verify this, let \(m\) be the midpoint of \([a,c]\). Midpoint
comparison bounds \(d(b,m)^2\) and \(d(d,m)^2\) by the corresponding
average squared distances to \(a,c\), minus \(d(a,c)^2/4\).
Combine these bounds with
\(d(b,d)^2\le2d(b,m)^2+2d(d,m)^2\) to obtain
\eqref{eq:cat0-quadrilateral}.

Approximate the four Euclidean domain points \(0,u,u+v,v\) by scaled
points of \(E_i-z\). Apply \eqref{eq:cat0-quadrilateral} to their images,
divide by \(t^2\), and use the two-point differential
\eqref{eq:two-point-metric-differential}. It follows that
\[
 N_z(u+v)^2+N_z(u-v)^2\le2N_z(u)^2+2N_z(v)^2.
\]
Apply the same inequality to \((u+v)/2\) and \((u-v)/2\), and use
homogeneity, to obtain the reverse inequality. The parallelogram law
therefore holds. Polarization produces a positive definite inner product
and the matrix \(G_i(z)\) in \eqref{eq:chart-metric}. Its entries are
measurable, since they are finite linear combinations of squared metric
differentials in fixed rational directions. Setting \(w'=0\) proves
\eqref{eq:centered-metric-differential}.
\end{proof}

The same bounds give, on each chart separately,
\begin{equation}
\label{eq:chart-metric-bounds}
 \Lip(\phi_i^{-1})^{-2}I\le G_i(z)\le\Lip(\phi_i)^2I,
 \qquad
 \Lip(\phi_i^{-1})^{-k}\le J_i(z)\le\Lip(\phi_i)^k.
\end{equation}
These constants may depend on the chart; no uniform bound across the
chart family will be needed.

For a Lipschitz scalar function \(u\) on \(Y\), define its chart covector
\(Du\) by differentiating \(u\circ\phi_i\) and transporting that covector
to \(Z_i\). Norms and inner products of these covectors use the dual
matrix \(G_i^{-1}\). Differentiating the scalar Lipschitz inequality at
density points gives
\begin{equation}
\label{eq:gradient-lipschitz}
 |Du|\le\Lip(u).
\end{equation}
Indeed, the derivative in every direction is bounded by
\(\Lip(u)N_z\), which is exactly the dual-norm bound. The chart covectors
obey the linear and product rules and the usual chain rule for continuously
differentiable scalar compositions. A scalar Lipschitz function has zero
chart derivative almost everywhere on any fixed level set: differentiate
at density points of that level set. These statements initially hold
almost everywhere in each Euclidean chart. Proposition~\ref{prop:chart-mass}
below identifies the same exceptional sets as mass-null sets.

\subsection{Quadratic mass bounds and exact chart mass}

The chart differential is Euclidean. The next two statements convert
that infinitesimal geometry into mass estimates with no dimensional
loss: a quadratic upper form controls mass from above, and inverse
chart coordinates recover the matching lower bound.

\begin{lemma}[A quadratic majorant for mass]
\label{lem:quadratic-mass}
Suppose that a finite-mass \(d\)-current \(C\) has a representation by
countably many determinant integrals with Lipschitz maps
\(\psi_i:F_i\to Y\), Borel sets \(F_i\subset\R^d\), and real signed
weights \(\vartheta_i\). Suppose that, at almost every point carrying
these weights, there is a measurable nonnegative quadratic form
\(P_i(z)\) such that
\begin{equation}
\label{eq:quadratic-distance-majorant}
 d(\psi_i(z+w),\psi_i(z))^2
 \le P_i(z)[w,w]+o(|w|^2),\qquad z+w\in F_i.
\end{equation}
If the measure
\begin{equation}
\label{eq:quadratic-mass-majorant}
 \eta=\sum_i(\psi_i)_\#
       \bigl(|\vartheta_i|\sqrt{\det P_i}\dd z\bigr)
\end{equation}
is finite, then \(\|C\|\le\eta\). No normality of the separate chart
currents is required.
\end{lemma}

\begin{proof}
Fix a tuple of \(d\) scalar \(1\)-Lipschitz tests. Almost everywhere on
each domain their compositions with \(\psi_i\) are differentiable in
the sense of scalar Lipschitz extensions. If \(\ell\) is one of the
resulting derivative rows, \eqref{eq:quadratic-distance-majorant} gives
\[
 |\ell(w)|\le\sqrt{P_i(z)[w,w]}\qquad(w\in\R^d).
\]
To see the bound in an arbitrary direction, approximate that direction
by rescaled domain points at a density-one point, and then use
differentiability and \eqref{eq:quadratic-distance-majorant}.
If \(P_i\) is positive definite, transform it to the identity and apply
Hadamard's inequality to the rows. Their absolute determinant is at most
\(\sqrt{\det P_i}\). If \(P_i\) is singular, all rows vanish on its
kernel; the full determinant and \(\det P_i\) both vanish, giving the
same bound. Substitution in the representation yields
\eqref{eq:mass-controlling-measure} with the measure \(\eta\).

The exceptional set is allowed to depend on the fixed test tuple. The
measure \(\eta\) is independent of that tuple and controls every tuple
separately; this is precisely what the definition of mass requires.
\end{proof}

\begin{proposition}[Exact chart mass]
\label{prop:chart-mass}
For the representation \eqref{eq:chart-action}, let \(G_i,J_i\) be as
in \eqref{eq:chart-metric}. Then
\begin{equation}
\label{eq:chart-mass}
 \|C\|=\sum_i(\phi_i)_\#\bigl(|\theta_i|J_i\dd z\bigr).
\end{equation}
Thus, with \(\rho_i=|\theta_i|J_i\), integration against the current
mass is integration against \(\rho_i\dd z\) on the disjoint charts.
\end{proposition}

\begin{proof}
Denote the right side of \eqref{eq:chart-mass} by \(\eta\), initially
allowing \(\eta(Y)=\infty\). We first prove \(\eta\le\|C\|\); this
will also establish finiteness before we apply the preceding lemma.

Fix \(0<\varepsilon<1\). Each chart domain, up to a Lebesgue-null set,
has a countable Borel partition into pieces \(E\) on which there is a
fixed positive definite matrix \(Q\), depending on the piece, satisfying
\begin{equation}
\label{eq:almost-isometric-chart-piece}
 (1-\varepsilon)|z-z'|_Q
 \le d(\phi_i(z),\phi_i(z'))
 \le(1+\varepsilon)|z-z'|_Q,
 \qquad z,z'\in E.
\end{equation}
Here is a construction. Approximate the measurable matrix \(G_i(z)\)
by an element of a countable dense family of positive definite matrices,
with a sufficiently small relative error. The centered estimate
\eqref{eq:centered-metric-differential} then gives a radius \(1/l\),
depending on the point, on which the two inequalities with this fixed
matrix and the prescribed \(\varepsilon\) hold against every other
point of the original chart domain. Group points according to the
matrix and \(l\), subdivide into sets of diameter less than \(1/l\),
and disjointify. These sets can be chosen Borel. Indeed, for any fixed
matrix and radius, the non-strict distance inequalities against all
domain points with \(0<|z-z'|<1/l\) can be tested on a fixed countable
dense subset of the domain, using continuity in \(z'\). All metric
differentiability points are covered by this countable construction.

At almost every density point of a piece \(E\), differentiation of
\eqref{eq:almost-isometric-chart-piece} in domain directions gives
\[
 |w|_{G_i(z)}\le(1+\varepsilon)|w|_Q,
 \qquad
 J_i(z)\le(1+\varepsilon)^k\sqrt{\det Q}.
\]
On \(\phi_i(E)\), every scalar coordinate of
\(Q^{1/2}\phi_i^{-1}\) is \((1-\varepsilon)^{-1}\)-Lipschitz by the
lower inequality in \eqref{eq:almost-isometric-chart-piece}.
By McShane's theorem \cite[Theorem~2.1]{Rieffel2006}, extend these
coordinates separately to all of \(Y\), preserving that Lipschitz
constant, and call the resulting tuple \(\pi\).

For any Borel subset \(D\subset E\), choose the bounded Borel first
coefficient supported on \(\phi_i(D)\) and equal there to
\(\sgn(\theta_i)\circ\phi_i^{-1}\). At density points of \(E\), the
extended coordinates have the same derivatives as \(Q^{1/2}z\), since
they agree with them on the piece. The images of all original charts
are disjoint. Hence \eqref{eq:chart-action}, applied to this coefficient
and tuple, gives exactly
\[
 \int_D|\theta_i|\sqrt{\det Q}\dd z
 = C(b,\pi_1,\ldots,\pi_k)
 \le(1-\varepsilon)^{-k}\|C\|(\phi_i(D)).
\]
The mass bound is legitimate for this Borel coefficient by its standard
extension from the first slot. Combining the last two estimates gives
\[
 \int_D|\theta_i|J_i\dd z
 \le\left(\frac{1+\varepsilon}{1-\varepsilon}\right)^k
       \|C\|(\phi_i(D)).
\]
For a Borel set \(B\subset Y\), take
\(D=E\cap\phi_i^{-1}(B)\) and sum over all pieces and all charts.
Their images are disjoint, while discarded Lebesgue-null sets make no
contribution to \(\eta\). Thus
\[
 \eta(B)\le
 \left(\frac{1+\varepsilon}{1-\varepsilon}\right)^k\|C\|(B).
\]
This proves that \(\eta\) is finite; letting
\(\varepsilon\downarrow0\) gives \(\eta\le\|C\|\).

Conversely, \eqref{eq:centered-metric-differential} gives the quadratic
majorant \(P_i=G_i\) for every chart. Its measure
\eqref{eq:quadratic-mass-majorant} is exactly the now finite measure
\(\eta\). Lemma~\ref{lem:quadratic-mass} yields \(\|C\|\le\eta\),
which completes the proof of the equality.
\end{proof}

We have established the exact mass formula, including arbitrary signed
integer multiplicities. Its next consequence is a lower density with
the Euclidean unit-ball coefficient; this is where integrality will
enter the sharp radial contradiction quantitatively.

\begin{lemma}[Euclidean lower density]
\label{lem:euclidean-density}
For an integral \(k\)-current \(C\) in \(Y\),
\begin{equation}
\label{eq:euclidean-lower-density}
 \liminf_{\rho\downarrow0}
       \rho^{-k}\|C\|(B(y,\rho))\ge\omega_k
 \qquad\text{for }\|C\|\text{-almost every }y.
\end{equation}
The balls in this statement are open.
\end{lemma}

\begin{proof}
By Proposition~\ref{prop:chart-mass}, it suffices to take
\(y=\phi_i(z)\), where \(z\) is a density-one point of \(E_i\), the
metric differential exists, \(z\) is a Lebesgue point of
\(\rho_i=|\theta_i|J_i\) extended by zero off \(E_i\), and
\(|\theta_i(z)|\ge1\). These conditions hold at almost every point
carrying chart mass. Fix \(\varepsilon>0\). For all sufficiently small
\(\rho\), the domain points in the ellipsoid
\[
 \{z+w:|w|_{G_i(z)}<\rho/(1+\varepsilon)\}
\]
map into \(B(y,\rho)\), by
\eqref{eq:centered-metric-differential}. This ellipsoid has volume
\[
 \frac{\omega_k\rho^k}{(1+\varepsilon)^kJ_i(z)}.
\]
Lebesgue differentiation on dilates of this fixed ellipsoid shows that
the contribution of this chart to \(\rho^{-k}\|C\|(B(y,\rho))\) has
lower limit at least
\(\omega_k|\theta_i(z)|/(1+\varepsilon)^k\).
Other charts contribute nonnegative mass. Letting
\(\varepsilon\downarrow0\) and using \(|\theta_i(z)|\ge1\) proves
\eqref{eq:euclidean-lower-density}.
\end{proof}

All subsequent chart covectors and their inner products are understood
with these metrics and the exact mass formula. In particular
\eqref{eq:gradient-lipschitz} holds \(\|C\|\)-almost everywhere.
For a cycle \(C\), the boundary action and the product rule give
\begin{equation}
\label{eq:cycle-product-rule}
 C(b,u,\pi_1,\ldots,\pi_{k-1})
 =-C(u,b,\pi_1,\ldots,\pi_{k-1})
\end{equation}
for Lipschitz scalar tests. Indeed, the sum of the two sides before
moving a term is \(C(1,bu,\pi)=\partial C(bu,\pi)=0\).
These identities concern the whole cycle. None of the preceding
arguments asserts that the separate restrictions to Borel chart images
have finite boundary mass.


The normalized cycle $A$ has all the chart formulas just proved. In the
next section we apply them to its radial deformation and to the swept
current, and compare the two masses through \eqref{eq:extremizer}.
